\documentclass[11pt]{article}
\usepackage{mathblatt}
\begin{document}
\weit
D: Leerzeile mit Strut nach jeder Reihe
\begin{gleichungsraster}[1]
\gl{x^2 = 4} & \gl{x^2 = 81} \\
\rule{0pt}{2mm} & \\
\gl{x^2 = 1} & \gl{x^2 = 100} \\
\rule{0pt}{2mm} & \\
\end{gleichungsraster}
Text danach D.

E: Box mit Tiefe am Zellende
\begin{gleichungsraster}[1]
\gl{x^2 = 4}\par\rule[-4mm]{0pt}{0pt} & \gl{x^2 = 81} \\
\gl{x^2 = 1}\par\rule[-4mm]{0pt}{0pt} & \gl{x^2 = 100} \\
\end{gleichungsraster}
Text danach E.
\end{document}
